\chapter{\textcolor{mycolor}{\textbf{Conception et architecture du système}}}
\markboth{\chaptername\ \thechapter.\ Conception et architecture du système}{}
\fancyhf{}
\fancyhead[L]{\small\itshape\nouppercase{\rightmark}}
\fancyfoot[L]{\small\itshape\nouppercase{\leftmark}}
\fancyfoot[R]{\small\itshape\thepage}
\renewcommand{\headrulewidth}{1pt}
\renewcommand{\footrulewidth}{0.4pt}

\vfill
\begin{flushright}{\Large\bfseries\textcolor{mycolor}{Introduction}}\end{flushright}
\addcontentsline{toc}{section}{Introduction}

Ce chapitre décrit la solution conçue pour le scénario pilote. Nous présentons d'abord l'architecture générale, la chaîne en sept étapes et les principes communs à toutes les sous-équipes. Nous détaillons ensuite la conception de chaque maillon : le matériel, le contrat d'interface et les données, la chaîne de perception, le moteur de décision et le backend, puis les deux applications.

\clearpage

\section{Architecture générale}

\subsection{La chaîne en sept étapes}

La solution est une chaîne de sept étapes, de la caméra jusqu'à l'historique (figure~\ref{fig:chaine}). Chaque
étape est tenue par une sous-équipe, et toutes se parlent par le même contrat d'interface.

\begin{figure}[H]
  \centering
  \begin{tikzpicture}[
      etape/.style={draw=mycolor, fill=mycolor!8, rounded corners=2pt, minimum width=3.1cm, minimum height=1.05cm,
                    align=center, font=\footnotesize},
      fleche/.style={-stealth, thick, mycolor}]
    \node[etape] (a) at (0,0) {Capteurs\\\scriptsize sous-équipe 1};
    \node[etape] (b) at (4.2,0) {Prétraitement\\\scriptsize sous-équipe 2};
    \node[etape] (c) at (8.4,0) {Analyse VLM\\\scriptsize sous-équipe 3};
    \node[etape] (d) at (8.4,-1.8) {Fusion contextuelle\\\scriptsize sous-équipe 4};
    \node[etape] (e) at (4.2,-1.8) {Décision SLM\\\scriptsize sous-équipe 4};
    \node[etape] (f) at (0,-1.8) {Notification mobile\\\scriptsize sous-équipe 6};
    \node[etape] (g) at (4.2,-3.6) {Historique local\\\scriptsize sous-équipes 2, 4, 5};
    \draw[fleche] (a) -- (b); \draw[fleche] (b) -- (c); \draw[fleche] (c) -- (d);
    \draw[fleche] (d) -- (e); \draw[fleche] (e) -- (f); \draw[fleche] (f) |- (g);
  \end{tikzpicture}
  \caption{La chaîne de traitement de LIFEWISE pour le scénario pilote}
  \label{fig:chaine}
\end{figure}

\subsection{Principes communs}

Quatre principes s'imposent à toutes les sous-équipes.

\begin{itemize}
  \item \textbf{Exécution locale.} Les traitements restent sur l'appareil ; aucune image n'est envoyée à un service distant.
  \item \textbf{Plusieurs sources d'image.} La caméra du Raspberry Pi, celle du téléphone ou une webcam peuvent alimenter la chaîne sans changer les modules. Le Raspberry Pi est le banc d'essai de l'Edge AI.
  \item \textbf{Un contrat d'interface commun.} Chaque étape reçoit et produit du JSON au format écrit par la sous-équipe Data ; un changement de format se propose et se relit.
  \item \textbf{Distinguer le vu de l'interprété.} La vision dit ce qu'elle voit, avec sa confiance ; la décision formule une invitation à vérifier, jamais un constat d'oubli.
\end{itemize}


% Sous-équipe 1 : IoT et matériel. Section du chapitre 3, appelée par rapport/main chapters.
% Retirer l'encadré \aRediger une fois la section écrite.

\section{Architecture matérielle et réseau}

\aRediger{1 (IoT et matériel)}{Décrire le choix de la caméra et des capteurs, leur positionnement près de la porte, le réseau et la liaison avec le téléphone.}

\subsection{Choix de la caméra}

\subsection{Positionnement de la caméra et des capteurs}

\subsection{Raspberry Pi et accès réseau}

\subsection{Liaison avec le téléphone}


% Sous-équipe 2 : Data. Section du chapitre 3, appelée par rapport/main chapters.
% Retirer l'encadré \aRediger une fois la section écrite.

\section{Contrat d'interface et modèle de données}

\aRediger{2 (Data)}{Décrire le contrat d'interface commun, l'objet contexte, le schéma SQLite et la politique de conservation.}

\subsection{Le contrat d'interface JSON}

\subsection{L'objet « contexte »}

\subsection{Le modèle de données SQLite}

\subsection{Politique de conservation et de consentement}


% Sous-équipe 3 : Vision (VLM). Section du chapitre 3, appelée par rapport/main chapters.
% Retirer l'encadré \aRediger une fois la section écrite.

\section{Conception de la chaîne de perception}

\aRediger{3 (Vision (VLM))}{Décrire la chaîne de traitement visuel, son format de sortie (objets et confiances) et la gestion des cas incertains.}

\subsection{Sélection des images}

\subsection{Entrées et sorties du module de vision}

\subsection{Gestion de la confiance et des observations incertaines}


% Sous-équipe 4 : Décision et backend. Section du chapitre 3, appelée par rapport/main chapters.
% Retirer l'encadré \aRediger une fois la section écrite.

\section{Conception du moteur de décision et du backend}

\aRediger{4 (Décision et backend)}{Décrire les règles, le rôle du SLM, l'architecture du service qui enchaîne les modules et l'API commune aux applications.}

\subsection{Règles de validation du scénario pilote}

\subsection{Rôle du petit modèle de langage}

\subsection{Architecture logicielle du service}

\subsection{Conception de l'API}


% Sous-équipe 5 : Application web. Section du chapitre 3, appelée par rapport/main chapters.
% Retirer l'encadré \aRediger une fois la section écrite.

\section{Conception du tableau de bord web}

\aRediger{5 (Application web)}{Décrire l'architecture du tableau de bord et ses écrans, dans la charte LIFEWISE.}

\subsection{Architecture web}

\subsection{Les écrans : historique, indicateurs, réglages}


% Sous-équipe 6 : Application mobile et coordination. Section du chapitre 3, appelée par rapport/main chapters.
% Retirer l'encadré \aRediger une fois la section écrite.

\section{Conception de l'application mobile}

\aRediger{6 (Application mobile et coordination)}{Décrire les écrans de l'application et le parcours d'une notification jusqu'à la réponse de l'utilisateur.}

\subsection{Les écrans de l'application}

\subsection{Réception et traitement de la notification}



\section*{Conclusion}
\addcontentsline{toc}{section}{Conclusion}

L'architecture est fixée et chaque maillon a ses entrées et ses sorties. Le chapitre suivant présente leur réalisation.
